\centerline{\bf \Huge Рыцари и лжецы. Рассуждения}
\title{ Рыцари и лжецы. Рассуждения}

\problem{1} На острове рыцарей и лжецов живут рыцари, которые всегда
говорят правду, и лжецы, которые всегда лгут. Однажды пятерых жителей этого острова по очереди спросили, сколько среди них рыцарей.\\
— Один, — ответил первый.\\
— Два, — ответил второй.\\
— Три, — ответил третий.\\
— Не верьте им, они все лжецы, — сказал четвёртый.\\
— Сам ты лжец! — сказал пятый четвёртому.\\
Сколько рыцарей было на самом деле?

\problem{2} На острове рыцарей и лжецов рыцари всегда говорят правду, а лжецы всегда лгут. В школе на этом острове учатся как рыцари, так и лжецы —
в одном классе. Однажды учитель спросил у четырёх детей — Ану, Бану, Вану и Дану — кто из них сделал домашнее задание. Они ответили:\\
— Ану: Домашнее задание сделали Бану, Вану и Дану.\\
— Бану: Домашнее задание не сделали Ану, Вану и Дану.\\
— Вану: Не верьте им, господин учитель! Ану и Бану — лжецы!\\
— Дану: Нет, господин учитель, Ану, Бану и Вану — рыцари!\\
Сколько рыцарей среди этих детей?

\problem{3} В самолёте летят жители города лжецов и жители города рыцарей. Рыцари всегда говорят правду, а лжецы всегда обманывают. Все пассажиры сели в ряды по 4 человека, и бортпроводник задал каждому пассажиру один и тот же вопрос. «Верно ли, что в вашем ряду столько же Ваших земляков, сколько жителей другого города?» Прозвучало ровно 70 утвердительных ответов. Сколько лжецов летит в самолёте? Человек считается своим собственным земляком.

\problem{4} В подводном царстве живут осьминоги с семью и восемью ногами.
Те, у кого 7 ног, всегда врут, а те, у кого 8 ног, всегда говорят правду. Однажды между тремя осьминогами состоялся такой разговор:\\
— Зелёный осьминог: У нас вместе 24 ноги.\\
— Синий осьминог: Ты прав!\\
— Красный осьминог: Глупости, Зелёный говорит ерунду!\\
Сколько ног было у каждого осьминога?\\

\problem{5} Страна имеет форму квадрата и разделена на 25 одинаковых квадратных графств. В каждом графстве правит либо граф-рыцарь, который всегда говорит правду, либо граф-лжец, который всегда лжет. Однажды каждый граф сказал: «Среди моих соседей поровну рыцарей и лжецов». Какое максимальное число рыцарей могло быть? (Графы являются соседями, если их графства имеют общую сторону.)

\problem{6} На острове рыцарей и лжецов живут рыцари, которые всегда говорят правду, и лжецы, которые всегда лгут. Однажды путешественник забрёл на вечеринку, на которой собрались 12 жителей острова (назовём их для краткости A, B, $\dotsc$ , L). Путешественник подсел к A и начал задавать вопросы: B $-$ рыцарь или лжец?; C — рыцарь или лжец?; $\dotsc$; L $-$ рыцарь или лжец? По полученным 11-ти ответам путешественник смог определить, сколько всего рыцарей среди A, $\dotsc$, L. Сделайте это и вы.

\problem{7} За круглым столом сидят 10 человек, каждый из которых либо рыцарь, который всегда говорит правду, либо лжец, который всегда лжёт. Двое из них заявили: «Оба моих соседа $-$ лжецы», а остальные восемь заявили: «Оба моих соседа — рыцари». Сколько рыцарей могло быть среди этих 10 человек?

\problem{8} На острове каждый житель либо рыцарь (всегда говорит правду), либо лжец (всегда лжет), либо обычный человек (может как говорить правду, так и лгать). Рыцари считаются людьми высшего ранга, обычные люди — среднего, а лжецы — низшего. А, В и – жители этого острова. Один из них — рыцарь, другой — лжец, а третий — обычный человек. А и В сказали следующее. А:
«В по рангу выше, чем С.» B: «С по рангу выше, чем А.» Что ответил С на вопрос: «Кто выше по рангу — А или В?»

\problem{9} На острове каждый житель либо рыцарь (всегда говорит правду), либо
лжец (всегда лжёт). Два жителя называются однотипными, если они либо оба рыцари, либо
оба лжецы. А, В и С — жители этого острова. А говорит: «В и С однотипны». Что ответит С на вопрос «А и В однотипны?»